% IU81001C
% Multiresistente Keime
% 14.0
% 2013-11-30
%
\label{m:mrsa-transport}\begin{itemize}
  \item Informationen über MRSA-Trägerschaft sowie Art/Ort der MRSA-Besiedelung
  \item Patientenvorbereitung und Transport
  \begin{itemize}
    \item Wunden müssen verbunden und gut abgedeckt sein
    \item Bei Besiedelung der Atemwege trägt der Patient einen Mund-/ Nasenschutz
    \item Vor dem Transport führt der Patient eine hygienische
    Händedesinfektion durch
  \end{itemize}
  \item Allgemeine Hygienemaßnahmen
  \begin{itemize}
    \item Einweghandschuhe und Schutzkittel
    \item Nach Ablegen sofortige hygienische Händedesinfektion
    \begin{itemize}
      \item Nach Abschluss des Patiententransportes sind alle Materialien,
      Geräte, Instrumente und Flächen, welche direkten Kontakt mit dem
      Patienten hatten, gemäß Hygieneplan zu
      desinfizieren.
      \item Alle waagerechten Oberflächen des Fahrzeuginnenraumes sind mit
      einem Flächendesinfektionsmittel in Form einer
      Scheuerwischdesinfektion desinfizierend zu reinigen.
      \item Abschließend führt das Personal eine hygienische
      Händedesinfektion durch.
      \item An Keimvehikel denken (Kugelschreiber, Haltegriffe, Stethoskop \ldots )!
    \end{itemize}
  \end{itemize}
  \item Nach Abschluss der vorangegangenen Hygiene- und
  Desinfektionsmaßnahmen ist das Fahrzeug wieder einsatzbereit.
\end{itemize}

Für die anderen Problemkeime \Z{(VRSA, VRE, ESBL,
{\dots})} gelten sinngemäß die gleichen
Verfahrensweisen.
}